\documentclass[11pt,a4paper]{article}

\usepackage[top=20mm,bottom=22mm,left=22mm,right=22mm]{geometry}
\usepackage{amsmath,amssymb}
\usepackage{booktabs}
\usepackage{enumitem}
\usepackage{titlesec}
\usepackage{hyperref}
\usepackage{microtype}
\usepackage{array}
\usepackage{setspace}
\usepackage{float}
\usepackage{caption}
\usepackage{graphicx}
\usepackage{subcaption}


\hypersetup{colorlinks=true,linkcolor=black,citecolor=black,urlcolor=black}

% ── Section: numbered, bold, small-caps feel, rule underneath ────────────────
\titleformat{\section}
  {\normalfont\normalsize\bfseries\uppercase}
  {\thesection.}{0.6em}{}
  [\vspace{-4pt}\rule{\linewidth}{0.4pt}\vspace{2pt}]
\titlespacing*{\section}{0pt}{10pt plus 2pt}{4pt plus 1pt}

% ── Subsection: bold, no number, compact ────────────────────────────────────
\titleformat{\subsection}[runin]
  {\normalfont\normalsize\bfseries}
  {}{0em}{}[.\enspace]
\titlespacing*{\subsection}{0pt}{6pt plus 1pt}{0pt}

\setlength{\parindent}{1.2em}
\setlength{\parskip}{3pt}

% ── List spacing ─────────────────────────────────────────────────────────────
\setlist{itemsep=1pt, topsep=3pt, parsep=0pt}

% =============================================================================
\begin{document}
\singlespacing

% ── Header block ─────────────────────────────────────────────────────────────
\begin{center}
  {\small Department of Mechanical Engineering,\enspace
   Indian Institute of Technology Madras,\enspace Chennai 600\,036}\\[4pt]
  \rule{\linewidth}{1.2pt}\\[6pt]
  {\large\bfseries Extended Abstract}\\[6pt]
  \rule{\linewidth}{0.6pt}
\end{center}

\vspace{4pt}

% ── Title and metadata ───────────────────────────────────────────────────────
\begin{center}
  {\normalsize\bfseries
   Autonomous Emergency Evasion at Blind Intersections:\\[2pt]
   From Formal Safety Theory to Hardware Deployment}
\end{center}

\vspace{6pt}

\noindent
\begin{tabular}{@{}p{0.18\linewidth}@{\hspace{4pt}}p{0.78\linewidth}@{}}
  \textbf{Student}  & Aadityanshu Abhinav \enspace (Roll No.\ ME22B088) \\[2pt]
  \textbf{Guide}    & Dr.\ Anuj Kumar Tiwari, Assistant Professor,
                      Department of Mechanical Engineering, IIT Madras \\
  % \textbf{Period}   & January -- May 2026 \\
\end{tabular}

\vspace{6pt}
\noindent\rule{\linewidth}{0.6pt}
\vspace{4pt}

% =============================================================================
\section{Introduction and Objectives}
% =============================================================================

Blind intersections -- junctions where cross-traffic is hidden by a building corner, a
parked truck, or a road curve -- represent one of the most dangerous situations in
autonomous driving. A vehicle approaching such a junction cannot determine whether a
threat exists until it is already within braking distance. Reactive planners that respond
only to observed obstacles are inadequate because the threat may appear too late to stop
for. Conversely, worst-case methods that assume a pursuer always exists force the vehicle
to stop far in advance of junctions that any human driver would pass without concern.

This project addresses the blind-intersection evasion problem in three stages:
(i)~implement and extend the game-theoretic occlusion-aware framework of Zhang and
Fisac~\cite{zhang2021occlusion} in the CARLA simulation environment and validate every
FSM state; (ii)~port the full planner stack to the PIX Moving Hooke autonomous
chassis equipped with a 32-beam 3D LiDAR, a 2D near-field guard, and a Jetson Orin
Nano; and (iii)~validate the complete evasion stack in outdoor hardware experiments at
3\,km/h with a Yujin Robots Kobuki mobile robot as the pursuer, demonstrating all four evasion outcomes -- push-through,
lateral dodge, curb climb, and reverse -- with zero contacts.

% =============================================================================
\section{Literature Review}
% =============================================================================

\subsection*{Forward reachability} The dominant prior approach computes a Forward
Reachable Set (FRS) for the worst-case hidden vehicle, requiring the ego to remain
outside the FRS or stop before it reaches the conflict zone. The FRS grows continuously
and does not shrink as the ego gains sensor evidence that parts of the occluded region
are empty. Neel and Saripalli~\cite{neel2020occluded} reduce this conservatism by
propagating velocity bounds derived from intersection geometry, but over-conservatism
remains inherent to the paradigm.

\subsection*{Game-theoretic planning} Zhang and Fisac~\cite{zhang2021occlusion} observe
that the over-conservatism of the FRS approach arises because it ignores the ego's
future ability to react once the pursuer becomes visible. They model the situation as a
two-mode zero-sum hybrid dynamic game: in the \emph{open-loop} phase (pursuer hidden),
the ego plays against the worst-case pursuer using a \emph{hidden set} -- the
forward reach-avoid set of all positions from which a pursuer could emerge; in the
\emph{closed-loop} phase (pursuer detected), the ego switches to a feedback evasion
strategy from the capture-basin reach-avoid game. This hybrid approach is provably safe
and demonstrably less conservative across three CARLA scenarios: a blind summit, a
T-intersection, and a truck-overtaking manoeuvre.

\subsection*{Gap addressed} None of the above works demonstrate hardware deployment on a
physical vehicle with a real sensor stack. This project bridges that gap end-to-end.

% =============================================================================
\section{Methodology}
% =============================================================================

\subsection*{Theoretical framework}
The hidden set $\tilde{H}(t)$ is maintained as a 2-D Shapely polygon and updated every
planning tick through four operations: (1)~union of newly occluded road area behind
LiDAR hits; (2)~outward propagation by $v_{\max,\text{pur}}\,\Delta t$;
(3)~subtraction of the confirmed-clear polygon from unblocked rays; and (4)~velocity-bound
pruning that prevents unbounded growth. The safe-passage condition (Theorem~1) requires
$\tilde{H}(t)\cap D(t)\approx\emptyset$, where $D(t)$ is the braking-distance corridor
ahead of the ego.

\subsection*{CARLA simulation}
Implemented in CARLA 0.9.16 (Town03) with a simulated 32-channel ray-cast LiDAR. Three
actors are used: a grey PIX Moving Hooke (ego), a red truck (occluder), and a red Tesla
Model 3 (pursuer). A five-state FSM (APPROACHING $\to$ PEEKING $\to$ PROCEEDING /
EVADING / YIELDING) governs the ego. A critical ego-yaw rotation correction was
identified and applied when converting sensor-local shadow angles to the world frame;
without it, the hidden set remains empty regardless of occlusion geometry.

\begin{figure}[H]
    \centering
    \includegraphics[width=0.82\linewidth]{CARLA_Scenario.png}
    \caption{A top view snapshot of the simulation scenario in CARLA: the grey vehicle travelling to the left is the agent, the red truck the occluder and the red car heading up is the pursuer.}
    \label{fig:carla_1}
\end{figure}

\begin{figure}[H]
    \centering
    \includegraphics[width=0.82\linewidth]{carla_iso.png}
    \caption{An isometric view snapshot of the simulation scenario in CARLA: the grey vehicle is the agent, while the red truck is the occluder; the state on the planner is \textbf{PEEKING} and the pursuer is hidden behind the occluder}
    \label{fig:carla_iso}
\end{figure}

\begin{figure}[htbp]
  \centering
  % First subfigure
  \begin{subfigure}[b]{0.45\textwidth}
    \includegraphics[width=\textwidth]{carla_purs1.png}
    \caption{Pursuer partially visible: planner in \textbf{PEEKING} state}
    \label{fig:carla_purs_vis1}
  \end{subfigure}
  \hfill
  % Second subfigure
  \begin{subfigure}[b]{0.45\textwidth}
    \includegraphics[width=\textwidth]{carla_purs_2.png}
    \caption{Pursuer fully visible: planner transitioning to \textbf{EVADING} state}
    \label{fig:carla_purs_vis2}
  \end{subfigure}
  
  \caption{Successive snapshots of the pursuer emerging from behind the occluder, invoking a transition from \textbf{PEEKING} to \textbf{EVADING} state}
  \label{fig:CARLA_trans}
\end{figure}

\subsection*{Hardware testbed}
The platform consists of: \textbf{Chassis} -- PIX Moving Hooke (2.6\,m\,$\times$\,1.7\,m),
CAN bus at 500\,kbit/s, 100\,ms VCU watchdog; \textbf{Primary sensor} -- Ouster OS1-32
(32-beam, 10\,Hz, 360$^\circ$\,H\,$\times$\,$\pm$22.5$^\circ$\,V), mounted upright at 1.3\,m above deck; \textbf{Secondary sensor} -- YDLidar X2 (2D, 8\,m),
front-mounted facing rearward, providing front and side collision guard coverage; \textbf{Compute} -- Jetson Orin
Nano 8\,GB, Ubuntu 22.04, Python 3.10, three concurrent processes using $<$55\,\% of
3 CPU cores and ${\sim}$340\,MB RAM. A Tata Altroz was used as the occluder, while a Yujin Robots Kobuki was used as the pursuer.

\subsection*{System architecture}
Three processes run concurrently. The \textbf{YDLidar Forwarder} reads the 2D LiDAR over USB and broadcasts a JSON safety packet (front, left, right distances) over UDP. The \textbf{Scenario Planner} reads the Ouster OS1-32 directly over Ethernet via a daemon thread, receives YDLidar packets, and runs the hidden-set tracker, FSM, motion planner, and all four LiDAR safety guards (Ouster front/side, YDLidar front/side) internally; guards override the FSM command by strict priority (estop $>$ brake $>$ drive). The \textbf{CAN Writer} transmits the resulting CAN frames at 20\,Hz. Only the CAN Writer writes to the vehicle bus; a crash of any single process does not leave the bus unattended.

\subsection*{Evasion priority stack}
When Theorem~1 is violated or the pursuer is detected, the planner evaluates four
options in order: (1)~\textbf{PUSH} -- sprint through if time-to-contact allows the ego
to clear first; (2)~\textbf{Lateral Dodge} -- bounded swerve of up to 0.8\,m toward the
curb if lane clearance permits; (3)~\textbf{Curb Climb} -- partial swerve using all
available clearance; (4)~\textbf{Reverse} -- back up along the approach path until the
OBB gap exceeds 5.0\,m. Near-field body separation is computed via the Separating Axis
Theorem (SAT) on oriented bounding boxes (OBBs).

\subsection*{Calibration and safety}
Steering encoder calibration (four setpoints, raw CAN value vs.\ hand-measured
road-wheel angle) confirmed decode divisor 22.0; consistent ${\sim}$2.5$^\circ$
under-read attributed to measurement parallax. Odometry closure error: $<$0.3\,m over a
16\,m square path. A human operator held the RC kill-switch -- the sole mechanism that
halts the vehicle independently of the script -- throughout every run. Operating speed
was capped at 0.833\,m/s (3\,km/h); all personnel outside a 5\,m exclusion zone. All
runs were conducted at night on deserted roads to avoid interaction with regular traffic;
no runs were conducted in rain due to LiDAR performance and electronics exposure
constraints.

% =============================================================================
\section{Summary of Work Completed}
% =============================================================================

\noindent\textbf{(a) Hidden-set tracker.} Complete update pipeline implemented in Python
(Shapely). Ego-yaw rotation bug identified and corrected. Runs at 10\,Hz on the Jetson;
each update $<$30\,ms (200-vertex polygon).

\noindent\textbf{(b) CARLA simulation.} FSM planner deployed and validated in CARLA
0.9.16. All three qualitative outcomes (PUSH, BRAKE, lateral dodge) observed across
repeated runs matching the baseline paper's Scenario~2 topology.

\noindent\textbf{(c) Hardware integration.} Ported planner to PIX Moving Hooke.
Implemented CAN odometry, pure-pursuit waypoint following, SAT-based OBB gap sensor,
pursuer velocity tracker, and three-process safety architecture. All calibration and
test scripts written and validated.

\noindent\textbf{(d) Hardware-specific additions.} PEEKING-HOLD condition added (hold
while pursuer tracker reports a converging cluster); all sensor filtering parameters
tuned for real outdoor noise conditions.

\noindent\textbf{(e) Outdoor experiments.} Twelve contested runs at 3\,km/h; results
summarised in Table~\ref{tab:results}. All four evasion outcomes demonstrated; zero
contacts.

\begin{table}[H]
\centering
\small
\begin{tabular}{lcc}
\toprule
Evasion outcome & Runs & Proportion \\
\midrule
PUSH (time-to-contact clear) & 3 & 25\,\% \\
BRAKE / EVADING              & 6 & 50\,\% \\
Lateral Dodge                & 2 & 17\,\% \\
Reverse (NR)                 & 1 & \phantom{0}8\,\% \\
\midrule
\textbf{Total contested runs} & \textbf{12} &  --  \\
\textbf{Contacts}             & \textbf{0}  &  --  \\
\bottomrule
\end{tabular}
\caption{\small Evasion outcome distribution across 12 hardware trials.}
\label{tab:results}
\end{table}

\noindent\textbf{(f) Key findings.}
\begin{itemize}
  \item The hidden-set framework correctly defers PROCEEDING when the shadow region is
        occupied, clearing within 3--5 ticks on an unoccluded approach.
  \item PEEKING-HOLD is necessary on hardware due to sensor noise; it conservatively
        delays PROCEEDING without introducing false EVADING transitions.
  \item At 3\,km/h the dominant outcome is BRAKE (${\sim}$0.08\,m braking distance always
        within Theorem~1's margin); PUSH and dodge are exercised only at tight timing
        windows, confirming the framework selects the correct action.
  \item Four simulation-to-hardware gaps were identified and bridged: speed cap
        (20\,km/h $\to$ 3\,km/h), pursuer type (vehicle $\to$ pedestrian), LiDAR
        noise, and odometry drift.
\end{itemize}

% ── References ───────────────────────────────────────────────────────────────
\renewcommand{\refname}{\normalsize\bfseries REFERENCES%
  \vspace{-4pt}\\\rule{\linewidth}{0.4pt}\vspace{2pt}}
{\small
\begin{thebibliography}{9}
\setlength{\itemsep}{1pt}

\bibitem{zhang2021occlusion}
Z.~Zhang and J.~F.~Fisac,
``Safe Occlusion-aware Autonomous Driving via Game-Theoretic Active Perception,''
\textit{arXiv:2105.08169}, 2021.

\bibitem{neel2020occluded}
G.~Neel and S.~Saripalli,
``Improving Bounds on Occluded Vehicle States for Use in Safe Motion Planning,''
in \textit{Proc.\ IEEE Int.\ Symp.\ Safety, Security, and Rescue Robotics (SSRR)},
pp.\ 268--275, 2020.

\bibitem{dosovitskiy2017carla}
A.~Dosovitskiy \textit{et al.},
``CARLA: An Open Urban Driving Simulator,''
in \textit{Proc.\ 1st Annual Conf.\ on Robot Learning}, pp.\ 1--16, 2017.

\end{thebibliography}
}

\end{document}
